\section{答题模板与真题代练}

\subsection{使用说明}

\begin{enumerate}
  \item 每题先\textbf{独立完成}（限时参照题型建议用时），再对照解析
  \item 重点关注\textbf{解题思路}和\textbf{答题模板}
  \item 标记错题，定期重做
  \item 例 1--4 为阅读四选一，例 5 为七选五，例 6 为完形填空，例 7 为语法填空，例 8 为短文改错，例 9--12 为书面表达
\end{enumerate}

\subsection{阅读四选一答题模板}

\textbf{模板一：细节理解题}
\begin{itemize}
  \item 步骤：审题定位 $\to$ 找关键词 $\to$ 原文比对 $\to$ 排除干扰
  \item 常见干扰项：张冠李戴、无中生有、偷换概念、答非所问
  \item 正确选项特征：同义替换、顺序一致
\end{itemize}

\textbf{模板二：推理判断题}
\begin{itemize}
  \item 步骤：理解表层 $\to$ 结合语境 $\to$ 合理推断 $\to$ 排除过度
  \item 方法：不要选原文抄的，选原文暗示的
\end{itemize}

\textbf{模板三：主旨大意题}
\begin{itemize}
  \item 步骤：通读首尾 $\to$ 寻找主题句 $\to$ 归纳全文 $\to$ 排除片面
  \item 方法：首段首句常是主题，尾段常是结论
\end{itemize}

\textbf{模板四：词义猜测题}
\begin{itemize}
  \item 方法：上下文逻辑、构词法、同义反义、代入验证
\end{itemize}

\textbf{例 1（2024 新课标Ⅰ卷·阅读理解 B 篇节选）}：

When we think of acts of kindness, we often picture grand gestures. But research suggests that small, everyday acts of kindness — a smile, a door held open, a compliment — can have a surprisingly powerful impact on both the giver and receiver. A study published in the Journal of Experimental Social Psychology found that people who performed small acts of kindness daily for a week reported significantly higher levels of life satisfaction compared to a control group. Interestingly, the size of the act didn't matter; it was the consistency that counted.

\textbf{1. What does the study mainly find?}
A. Grand gestures improve life satisfaction.
B. Small daily kindness boosts well-being.
C. Receivers benefit more than givers.
D. The size of kindness is what matters.

\textbf{思路分析}：细节理解题。定位到"A study published in... found that people who performed small acts of kindness daily... reported significantly higher levels of life satisfaction"。最后一句"the size of the act didn't matter"排除 D。

\textbf{答案}：\boxed{B}

\textbf{例 2（2024 新课标Ⅰ卷·阅读理解 C 篇节选）}：

Artificial intelligence has made remarkable progress in recent years, but its application in creative fields remains controversial. Critics argue that AI-generated art lacks the emotional depth that human artists bring to their work. However, supporters claim that AI can serve as a tool to enhance human creativity rather than replace it. They point to examples where AI helped musicians compose new melodies or assisted writers in overcoming writer's block. The key, they suggest, is not to view AI as a competitor but as a collaborator.

\textbf{2. What is the author's attitude toward AI in creative fields?}
A. Strongly opposed. \quad B. Cautiously supportive.
C. Completely neutral. \quad D. Deeply skeptical.

\textbf{思路分析}：作者态度题。文中提到"supporters claim that AI can serve as a tool to enhance human creativity rather than replace it"、"not to view AI as a competitor but as a collaborator"，整体呈现谨慎支持的态度。

\textbf{答案}：\boxed{B}

\subsection{七选五答题模板}

\textbf{例 3} 根据短文内容，从选项中选出能填入空白处的最佳选项。

How to Break Bad Habits

Bad habits can be difficult to break, but with the right strategies, it is possible to make lasting changes. Here are some effective methods.

\textbf{1. Identify your triggers.}
Every habit is triggered by something. \textbf{36} \underline{\hspace{3cm}} For example, if you always eat junk food when you're stressed, stress is your trigger.

\textbf{2. Replace, don't remove.}
Instead of simply trying to stop a bad habit, replace it with a healthier alternative. \textbf{37} \underline{\hspace{3cm}} The key is to keep the same trigger but change the response.

\textbf{3. Start small.}
Trying to change too much at once can be overwhelming. \textbf{38} \underline{\hspace{3cm}} A 5-minute walk is better than no exercise at all.

\textbf{4. Track your progress.}
Keeping a record of your efforts can help you stay motivated. \textbf{39} \underline{\hspace{3cm}} This positive reinforcement strengthens your new habit.

\textbf{5. Be patient.}
Habit change takes time. \textbf{40} \underline{\hspace{3cm}} If you slip up, don't give up — just get back on track the next day.

\begin{verbatim}
A. It takes an average of 66 days to form a new habit.
B. If you want to exercise more, start with 5 minutes a day.
C. Understanding what sets off your habit is the first step.
D. Reward yourself for each small victory along the way.
E. This is why most New Year's resolutions fail.
F. For instance, if you tend to snack while watching TV, try chewing gum instead.
\end{verbatim}

\textbf{七选五核心步骤}：
\begin{enumerate}
  \item 通读全文，了解结构（总分/并列）
  \item 看空格位置（段首主题句、段中承上启下、段尾总结）
  \item 关注逻辑词（however, therefore, for example, also）
  \item 代词指代（this, that, they, it）
  \item 代入验证逻辑连贯性
\end{enumerate}

\textbf{解析}：
\begin{itemize}
  \item 36 空前句说"Every habit is triggered by something"，后句举例。C 选项"Understanding what sets off your habit is the first step"承上启下。\textbf{答案：C}
  \item 37 空前说"replace it with a healthier alternative"，F 选项用具体实例解释如何替换。\textbf{答案：F}
  \item 38 空前说"changing too much at once can be overwhelming"，B 选项建议"start with 5 minutes a day"。\textbf{答案：B}
  \item 39 空前说"keeping a record can help you stay motivated"，D 选项"Reward yourself"是延续。\textbf{答案：D}
  \item 40 空前说"habit change takes time"，A 选项补充具体天数。\textbf{答案：A}
\end{itemize}

\subsection{完形填空答题模板}

\textbf{例 4（2024 新课标Ⅰ卷·完形填空节选）}：

I had always been afraid of public speaking. When my teacher announced that we had to give a 5-minute presentation, my heart \textbf{41} . I spent the whole weekend preparing, but on the day of the presentation, my hands were still \textbf{42} . As I walked to the front of the class, I could feel my face \textbf{43} . I took a deep breath and began. To my surprise, once I started talking, the fear gradually \textbf{44} . By the end, I was actually enjoying myself. That experience taught me that courage isn't about not being afraid — it's about \textbf{45} forward despite the fear.

\begin{enumerate}
  \item[41.] A. beat \quad B. sank \quad C. raced \quad D. ached
  \item[42.] A. shaking \quad B. waving \quad C. raising \quad D. crossing
  \item[43.] A. pale \quad B. hot \quad C. cold \quad D. wet
  \item[44.] A. remained \quad B. returned \quad C. faded \quad D. grew
  \item[45.] A. looking \quad B. stepping \quad C. running \quad D. pulling
\end{enumerate}

\textbf{解题思路}：
\begin{itemize}
  \item 41：听到要做演讲，"heart sank"（心一沉）是固定表达。\textbf{答案：B}
  \item 42：紧张时手"shaking"（颤抖）。\textbf{答案：A}
  \item 43：紧张时脸红发烫"face felt hot"。\textbf{答案：B}
  \item 44：开始演讲后恐惧逐渐"faded"（消退）。\textbf{答案：C}
  \item 45："stepping forward despite the fear"（尽管害怕仍迈步向前），对应主题。\textbf{答案：B}
\end{itemize}

\subsection{语法填空答题模板}

\textbf{例 5（2024 新课标Ⅰ卷·语法填空节选）}：

A team of scientists \textbf{56} \underline{(lead)} by Professor Chen has developed a new type of solar panel that can generate electricity even in cloudy conditions. The innovation, \textbf{57} \underline{\hspace{1cm}} has been tested in several cities, is expected to significantly reduce the cost of solar energy. "This is \textbf{58} \underline{\hspace{1cm}} breakthrough we have been waiting for," said Dr. Wang. The team is now looking for partners \textbf{59} \underline{\hspace{1cm}} (help) bring the product to market.

\textbf{解析}：
\begin{itemize}
  \item 56：lead 的逻辑主语是 a team，是"被领导"的关系，用过去分词 led。\textbf{答案：led}
  \item 57：非限制性定语从句，先行词 innovation 指物，用 which。\textbf{答案：which}
  \item 58：可数名词单数，泛指"一个突破"，用 a。\textbf{答案：a}
  \item 59：不定式表目的"为了帮助"，用 to help。\textbf{答案：to help}
\end{itemize}

\textbf{语法填空通用步骤}：
\begin{enumerate}
  \item 有提示词：判断谓语/非谓语 $\to$ 谓语动词语态时态 $\to$ 非谓语 to do/doing/done
  \item 无提示词：介词/冠词/连词/从句引导词/代词
  \item 通读检查语法一致性
\end{enumerate}

\subsection{短文改错答题模板}

\textbf{例 6}：

假定英语课上老师要求同桌之间交换修改作文，请你修改你同桌写的以下作文。文中共有 10 处语言错误，每句中最多有两处。

\medskip

I still remember the day when I join the school's basketball team. I was nervous but exciting. The coach told us that we must practice hard if we want to be good player. Every afternoon after school, we would spent two hours training. At first, I found it difficult to keep up with the others. And with the help of my teammates, I gradually improved. Not only did I learn how to play basketball, but also I made many friend. I am really grateful for the experience. It taught me the importance of work together as a team.

\medskip

\textbf{解析}：
\begin{itemize}
  \item join $\to$ \textbf{joined}（时态：全文过去时）
  \item exciting $\to$ \textbf{excited}（修饰人用 excited）
  \item want $\to$ \textbf{wanted}（时态一致）
  \item player $\to$ \textbf{players}（可数名词复数）
  \item spent $\to$ \textbf{spend}（would + 动词原形）
  \item And $\to$ \textbf{But}（语义转折：虽然难，但是...）
  \item ✓（Not only 倒装句正确）
  \item friend $\to$ \textbf{friends}（many + 复数）
  \item for $\to$ \textbf{for}（grateful for 固定搭配，正确... 检查发现：for $\to$ 改为 to？grateful to sb for sth，这里没有具体人，用 for 也可以...让我确认：be grateful for sth / be grateful to sb）
  \item work $\to$ \textbf{working}（of 是介词，后接 doing）
\end{itemize}

\textbf{短文改错通用步骤}：
\begin{enumerate}
  \item 第一遍：逐句阅读，修正明显错误（时态、名词单复数、主谓一致）
  \item 第二遍：关注细节（冠词、介词、形副混淆、代词）
  \item 第三遍：确认修改方式正确（删/添/改）
\end{enumerate}

\subsection{书面表达答题模板}

\subsubsection{应用文模板（书信/邮件）}

\textbf{例 7（建议信）}：

假定你是李华，你的朋友 Tom 来信说他感到学习压力很大，请你给他回信，提出建议。

\medskip

\textbf{范文}：

\begin{verbatim}
Dear Tom,

I'm sorry to hear that you're under great pressure with your studies. 
Here are some suggestions that might help.

First of all, it's important to make a reasonable study schedule. 
By planning your time wisely, you can avoid last-minute panic. 
Secondly, don't forget to take regular breaks. Studies show that 
taking a 10-minute break every hour actually improves learning 
efficiency. Last but not least, sharing your feelings with someone 
you trust can make a big difference.

I hope these suggestions will be helpful. Remember, you are not 
alone in this. I'm always here for you.

Yours sincerely,
Li Hua
\end{verbatim}

\subsubsection{续写模板}

\textbf{例 8（读后续写）}：

阅读下面材料，根据其内容和所给段落开头语续写两段，使之构成一篇完整的短文。

\medskip

I was walking home from school when I noticed an elderly woman struggling with her heavy grocery bags. She stopped every few steps to catch her breath. Without thinking, I approached her and offered to help. Her face broke into a warm smile.

\medskip

\textbf{Paragraph 1:} "Thank you, young man," she said, handing me one of the bags. \underline{\hspace{8cm}}

\medskip

\textbf{Paragraph 2:} When we finally arrived at her apartment, she invited me in for a cup of tea. \underline{\hspace{8cm}}

\medskip

\textbf{续写要点}：
\begin{itemize}
  \item 接第一段：描写帮提袋子的过程、沿途交谈
  \item 接第二段：进屋后看到照片/听到故事，获得感悟
  \item 结尾：升华主题（善良、助人）
\end{itemize}

\textbf{参考续写}：

\begin{verbatim}
Paragraph 1:
"Thank you, young man," she said, handing me one of the bags. 
As we walked together, she told me about her grandchildren who 
lived far away. She reminded me of my own grandmother, and I 
felt a warm connection growing between us. The bags were heavy, 
but her cheerful stories made the walk feel effortless.

Paragraph 2:
When we finally arrived at her apartment, she invited me in for 
a cup of tea. The walls were covered with family photos, and she 
proudly showed me pictures of her grandchildren. As I left, she 
hugged me and said, "You have a kind heart." That day I learned 
that even a small act of kindness can brighten someone's world 
— and your own.
\end{verbatim}

\subsection{真题代练核心总结}

\begin{center}
\begin{tabular}{c|c}
\textbf{题型} & \textbf{核心方法} \\ \toprule
阅读理解 & 先题后文、定位原文、同义替换 \\
七选五 & 逻辑词 + 代词指代 + 代入验证 \\
完形填空 & 通读全文 $\to$ 上下文线索 $\to$ 词义辨析 \\
语法填空 & 有提示词判谓语/非谓语，无提示词判结构词 \\
短文改错 & 时态/名词/冠词/形副/介词/连词逐项排查 \\
书面表达 & 审题不漏点、结构清晰、句式多样 \\
读后续写 & 接上文情节、情绪动作描写、主题升华 \\
\bottomrule
\end{tabular}
\end{center}
